\documentclass[11pt]{article}
\usepackage[margin=1.1in]{geometry}
\usepackage{amsmath,amssymb,amsthm,mathtools,booktabs,array,tabularx,longtable}
\usepackage[dvipsnames]{xcolor}
\usepackage[colorlinks=true,linkcolor=MidnightBlue,citecolor=MidnightBlue,urlcolor=MidnightBlue]{hyperref}
\usepackage{fancyhdr}
\pagestyle{fancy}\fancyhf{}\renewcommand{\headrulewidth}{0pt}
\fancyfoot[R]{\thepage}
\fancyfoot[L]{\scriptsize\color{gray}Generated by the Transport Lab from approved parts of its store; not a deposited paper. Papers of record: \href{https://zenodo.org/communities/transport-papers}{Zenodo, Transport Papers}}
\fancypagestyle{plain}{\fancyhf{}\fancyfoot[R]{\thepage}\fancyfoot[L]{\scriptsize\color{gray}Generated by the Transport Lab from approved parts of its store; not a deposited paper. Papers of record: \href{https://zenodo.org/communities/transport-papers}{Zenodo, Transport Papers}}\renewcommand{\headrulewidth}{0pt}}
\newtheorem{theorem}{Theorem}[section]
\newtheorem{proposition}[theorem]{Proposition}
\newtheorem{lemma}[theorem]{Lemma}
\newtheorem{corollary}[theorem]{Corollary}
\newtheorem{conjecture}[theorem]{Conjecture}
\newtheorem{refutation}[theorem]{Refutation}
\newtheorem{prediction}[theorem]{Prediction}
\newtheorem{principle}[theorem]{Principle}
\newtheorem{observation}[theorem]{Observation}
\newtheorem{axiom}[theorem]{Axiom}
\theoremstyle{definition}
\newtheorem{definition}[theorem]{Definition}
\newtheorem{fact}[theorem]{Fact}
\newtheorem{assembled}[theorem]{Assembly}
\newtheorem{postulate}[theorem]{Postulate}
\newtheorem{premise}[theorem]{Premise}
\newtheorem{hypothesis}[theorem]{Hypothesis}
\newtheorem{remark}[theorem]{Remark}
\newtheorem{example}[theorem]{Example}
\theoremstyle{plain}
\setlength{\marginparwidth}{0.8in}
\newcommand{\changedmark}[2]{\marginpar{\raggedright\scriptsize\color{#1}#2}}
\providecommand{\Aa}{}\renewcommand{\Aa}{\mathcal{A}}
\providecommand{\Ad}{}\renewcommand{\Ad}{\widetilde{\mathrm{Ad}}}
\providecommand{\Aev}{}\renewcommand{\Aev}{\mathcal{A}^{+}}
\providecommand{\Aod}{}\renewcommand{\Aod}{\mathcal{A}^{-}}
\providecommand{\B}{}\renewcommand{\B}{\mathrm{B}}
\providecommand{\C}{}\renewcommand{\C}{\mathbb{C}}
\providecommand{\CAR}{}\renewcommand{\CAR}{\mathrm{CAR}}
\providecommand{\CS}{}\renewcommand{\CS}{{\mathbb{S}^{+}}}
\providecommand{\CSm}{}\renewcommand{\CSm}{{\mathbb{S}^{-}}}
\providecommand{\Cl}{}\renewcommand{\Cl}{\mathrm{Cl}}
\providecommand{\Con}{}\renewcommand{\Con}{\mathcal{C}}
\providecommand{\Dl}{}\renewcommand{\Dl}{\Delta}
\providecommand{\EN}{}\renewcommand{\EN}{E_{N}}
\providecommand{\End}{}\renewcommand{\End}{\mathrm{End}}
\providecommand{\F}{}\renewcommand{\F}{\mathbb{F}}
\providecommand{\G}{}\renewcommand{\G}{\Gamma}
\providecommand{\HH}{}\renewcommand{\HH}{\mathbb{H}}
\providecommand{\IS}{}\renewcommand{\IS}{I\sigma_3}
\providecommand{\J}{}\renewcommand{\J}{\mathbf{J}}
\providecommand{\LH}{}\renewcommand{\LH}{L_{H}}
\providecommand{\Mat}{}\renewcommand{\Mat}{\mathrm{M}}
\providecommand{\Mo}{}\renewcommand{\Mo}{M^{\circ}}
\providecommand{\Mp}{}\renewcommand{\Mp}{M^{+}}
\providecommand{\Msun}{}\renewcommand{\Msun}{M_{\odot}}
\providecommand{\Mt}{}\renewcommand{\Mt}{\widetilde{M}}
\providecommand{\Ncone}{}\renewcommand{\Ncone}{\mathcal{N}}
\providecommand{\Neff}{}\renewcommand{\Neff}{N_{\mathrm{eff}}}
\providecommand{\PCS}{}\renewcommand{\PCS}{\mathbb{P}{\mathbb{S}^{+}}}
\providecommand{\PT}{}\renewcommand{\PT}{\mathbb{PT}}
\providecommand{\Pin}{}\renewcommand{\Pin}{\mathrm{Pin}}
\providecommand{\Q}{}\renewcommand{\Q}{\mathbb{Q}}
\providecommand{\R}{}\renewcommand{\R}{\mathbb{R}}
\providecommand{\Ric}{}\renewcommand{\Ric}{\mathrm{Ric}}
\providecommand{\SBH}{}\renewcommand{\SBH}{S_{\mathrm{BH}}}
\providecommand{\Sc}{}\renewcommand{\Sc}{\mathbb{S}}
\providecommand{\Sch}{}\renewcommand{\Sch}{\mathrm{P}}
\providecommand{\Sig}{}\renewcommand{\Sig}{\Sigma}
\providecommand{\Som}{}\renewcommand{\Som}{S_\omega}
\providecommand{\Sp}{}\renewcommand{\Sp}{\mathbb{S}}
\providecommand{\Sph}{}\renewcommand{\Sph}{\mathrm{S}}
\providecommand{\Spi}{}\renewcommand{\Spi}{S_\pi}
\providecommand{\Spin}{}\renewcommand{\Spin}{\mathrm{Spin}}
\providecommand{\Sth}{}\renewcommand{\Sth}{S_{\mathrm{th}}}
\providecommand{\T}{}\renewcommand{\T}{\mathbb{T}}
\providecommand{\TH}{}\renewcommand{\TH}{T_{\mathrm{H}}}
\providecommand{\Tr}{}\renewcommand{\Tr}{\operatorname{Tr}}
\providecommand{\UHF}{}\renewcommand{\UHF}{\mathrm{UHF}}
\providecommand{\Z}{}\renewcommand{\Z}{\mathbb{Z}}
\providecommand{\Zc}{}\renewcommand{\Zc}{\mathrm{Z}}
\providecommand{\cmg}{}\renewcommand{\cmg}{\mathrm{cm^{2}\,g^{-1}}}
\providecommand{\dS}{}\renewcommand{\dS}{\mathrm{dS}}
\providecommand{\dsc}{}\renewcommand{\dsc}{\mathrm{dsc}}
\providecommand{\dtPG}{}\renewcommand{\dtPG}{dt}
\providecommand{\dtS}{}\renewcommand{\dtS}{dt_{\mathrm{S}}}
\providecommand{\ellv}{}\renewcommand{\ellv}{\ell}
\providecommand{\er}{}\renewcommand{\er}{e_r}
\providecommand{\et}{}\renewcommand{\et}{e_t}
\providecommand{\ev}{}\renewcommand{\ev}{\mathrm{ev}}
\providecommand{\gl}{}\renewcommand{\gl}{\mathfrak{gl}}
\providecommand{\hatV}{}\renewcommand{\hatV}{\widehat{V}}
\providecommand{\hb}{}\renewcommand{\hb}{\bar{h}}
\providecommand{\hf}{}\renewcommand{\hf}{\hat{f}}
\providecommand{\id}{}\renewcommand{\id}{\mathrm{id}}
\providecommand{\kpc}{}\renewcommand{\kpc}{\mathrm{kpc}}
\providecommand{\lP}{}\renewcommand{\lP}{\ell_{P}}
\providecommand{\mD}{}\renewcommand{\mD}{\mathbf{m}_{D}}
\providecommand{\mb}{}\renewcommand{\mb}{\bar m}
\providecommand{\mnu}{}\renewcommand{\mnu}{\Sigma m_{\nu}}
\providecommand{\nInf}{}\renewcommand{\nInf}{n_\infty}
\providecommand{\nO}{}\renewcommand{\nO}{n_0}
\providecommand{\ninf}{}\renewcommand{\ninf}{n_{\infty}}
\providecommand{\nz}{}\renewcommand{\nz}{n_{0}}
\providecommand{\proj}{}\renewcommand{\proj}{\operatorname{proj}}
\providecommand{\pv}{}\renewcommand{\pv}{\pi}
\providecommand{\rad}{}\renewcommand{\rad}{\mathrm{rad}}
\providecommand{\sgn}{}\renewcommand{\sgn}{\operatorname{sgn}}
\providecommand{\shA}{}\renewcommand{\shA}{\mathcal{A}}
\providecommand{\shC}{}\renewcommand{\shC}{\mathcal{C}}
\providecommand{\shE}{}\renewcommand{\shE}{\mathcal{E}}
\providecommand{\shJ}{}\renewcommand{\shJ}{\mathcal{J}}
\providecommand{\shK}{}\renewcommand{\shK}{\mathcal{K}}
\providecommand{\shL}{}\renewcommand{\shL}{L}
\providecommand{\shW}{}\renewcommand{\shW}{\mathcal{W}}
\providecommand{\sll}{}\renewcommand{\sll}{\mathfrak{sl}}
\providecommand{\so}{}\renewcommand{\so}{\mathfrak{so}}
\providecommand{\su}{}\renewcommand{\su}{\mathfrak{su}}
\providecommand{\supp}{}\renewcommand{\supp}{\operatorname{supp}}
\providecommand{\tr}{}\renewcommand{\tr}{\operatorname{tr}}
\providecommand{\uu}{}\renewcommand{\uu}{\mathfrak{u}}
\providecommand{\br}{}\renewcommand{\br}[1]{\langle #1\rangle}
\providecommand{\bra}{}\renewcommand{\bra}[1]{\langle #1|}
\providecommand{\file}{}\renewcommand{\file}[1]{\texttt{#1}}
\providecommand{\hp}{}\renewcommand{\hp}[2]{( #1 \,|\, #2 )}
\providecommand{\ip}{}\renewcommand{\ip}[2]{\langle #1, #2\rangle}
\providecommand{\ket}{}\renewcommand{\ket}[1]{|#1\rangle}
\providecommand{\rel}{}\renewcommand{\rel}[2]{D\bigl(#1 \,\|\, #2\bigr)}
\providecommand{\rev}{}\renewcommand{\rev}[1]{\widetilde{#1}}
\providecommand{\slashed}{}\renewcommand{\slashed}[1]{{#1}\mkern-10mu/}
\title{\bfseries The singularity}
\author{Leon Butler\\ \normalsize\href{mailto:lbutler.math@gmail.com}{\texttt{lbutler.math@gmail.com}}\\ \normalsize ORCID: \href{https://orcid.org/0009-0006-4449-5617}{0009-0006-4449-5617}}
\date{Generated by the Transport Lab, 2 November 2026}
\begin{document}
\maketitle
\begin{abstract}
\normalsize
This trail goes to the centre of a non-rotating black hole and asks exactly what fails there. The gauge field describing the black hole is the identity plus a shear $P$ with $P^2 = 0$, so it can be inverted at every radius, the event horizon included.

All of the blow-up lives instead in the connection and the curvature. A curvature invariant that no gauge choice can alter grows without bound at the centre, so the breakdown is real rather than an artefact.

Once the growth is divided out, the field settles to a degenerate limit with a null direction that has no partner: an obstruction of type N. Adding a hyperbolic plane, coupled to the lost direction only inside the horizon, repairs the geometry everywhere, the centre included.
\end{abstract}
\newpage
\noindent \textbf{Authorship and disclosure. }The founding idea of the series is the author's: that division by zero in a Clifford algebra should be read not as a prohibition but as a transport into a higher-dimensional algebra, and that the transported division need not be accompanied by an inverse operation. The direction of investigation and the critical review of the results are likewise the author's. The mathematical development and the drafting of the text were carried out in extended collaboration with Claude Fable 5 (Anthropic), an AI system.

\noindent\fbox{\parbox{\dimexpr\linewidth-2\fboxsep-2\fboxrule}{\textbf{This is not a deposited paper.} It was generated by the Transport Lab on 2 November 2026 and assembled from approved parts of its store. It follows results across several papers of the series; each is stated with its proof and credited to the paper that first proved it. The official versions of those papers are on Zenodo, in the Transport Papers community, \url{https://zenodo.org/communities/transport-papers}.}}
\newpage

\section*{Introduction}
The walk follows the order of the argument: first the gauge field and its connection, then where the divergence sits and why it cannot be removed, then the exact shape of the degeneration and what the series makes of it.

\section{The singularity}

Each step below narrows down where, and in what sense, the black-hole solution breaks at its centre.

We start with the gravitational field itself, which in Gauge Theory Gravity is a linear map $\hb$ acting on the directions of flat spacetime.

In the Painlev\'e-Gullstrand gauge it is $\hb = 1 - \beta P$ with $P^2 = 0$ exactly, so its inverse is \[ \hb\^{}{-1} = 1 + \beta P \] at every $r > 0$. At the horizon, $\beta = 1$, division does not fail: the horizon is algebraically unremarkable.

\begin{proposition}[The unipotent ray]\label{c:singularity.unipotent-ray}
In this formulation of gravity, the field describing a non-rotating black hole is the identity plus a shear whose square is zero. That makes it reversible at every radius, with an explicit formula for the inverse. In particular nothing goes wrong algebraically at the event horizon.
\emph{First proved in} T5, Prop 2.2 (\href{https://doi.org/10.5281/zenodo.22721790}{doi:10.5281/zenodo.22721790}).
\end{proposition}

The field $\hb$ is only half of the description; the other half is the rotation-gauge field $\Omega$, which records how local frames boost and rotate from point to point.

Here it is a multiple of the single radial boost bivector $\sigma_r$: \[ \Omega(\partial\_t) = \frac{M}{r\^{}2}\,\sigma\_r, \qquad \Omega(\partial\_r) = \sqrt{\frac{M}{2r\^{}3}}\;\sigma\_r \] Along time it is Newton's field strength, while along the radius it grows without bound towards the centre.

\begin{proposition}[The connection bivector, computed]\label{c:singularity.connection-bivector-computed}
In the gauge-theory description of gravity, the connection records how local frames boost and rotate from point to point. For a non-rotating black hole in this gauge it can be written down exactly. Along the time direction it equals Newton's gravitational field strength, mass over radius squared, at every radius; along the radial direction it grows without bound towards the centre, at the square root of the rate at which the field strength grows.
\emph{First proved in} T5, Prop 2.3 (\href{https://doi.org/10.5281/zenodo.22721790}{doi:10.5281/zenodo.22721790}).
\end{proposition}

Putting the two halves together locates the trouble.

The pair solves the vacuum equations at every $r > 0$ with $\det \hb = 1$ throughout, so the divergence lives in $\Omega$ and in the curvature, not in $\hb$. Reading this as a statement about Einstein's theory uses the imported equivalence of the two, wherever $\hb$ is invertible.

\begin{proposition}[The solution, and where the divergence lies]\label{c:singularity.solution-where-divergence-lies}
The gravitational field of a non-rotating black hole is written here in a gauge-theory language with two ingredients: a field describing how lengths are stretched, and a connection describing how directions turn. The first ingredient stays well behaved at every distance from the centre and never squeezes or expands volume. All of the blow-up as the centre is approached lives in the connection, in the curvature built from it, and in the ordinary metric components.
\emph{First proved in} T5, Prop 2.4 (\href{https://doi.org/10.5281/zenodo.22721790}{doi:10.5281/zenodo.22721790}).
\end{proposition}
\noindent\emph{Rests on:} \ref{c:singularity.connection-bivector-computed}, \ref{c:singularity.imported-equivalence-einstein-equations}.

Since the curvature carries the divergence, it is worth seeing its form exactly.

Acting on bivectors $B$, the field strength is a fixed shape times one scalar, \[ $\mathcal{R}$(B) = -\frac{M}{r\^{}3}\,S(B), \qquad S = 1 - 3\Pi \] where $\Pi$ projects onto the radial plane and its dual. The shape $S$, with spectrum $(-2,-2,1,1,1,1)$, never changes; only $M/r^3$ grows.

\begin{theorem}[The curvature of the gauge field: exact shape decoupling]\label{c:singularity.curvature-gauge-field-exact-shape-decoupling}
The tidal curvature of a non-rotating black hole splits exactly into a fixed geometric shape times a single number. The shape singles out the radial direction and is the same at every radius; only the strength, mass divided by radius cubed, changes. So the whole approach to the central singularity is governed by that one scalar.
\emph{First proved in} T5, Thm 6.1 (\href{https://doi.org/10.5281/zenodo.22721790}{doi:10.5281/zenodo.22721790}).
\end{theorem}
\noindent\emph{Rests on:} \ref{c:pairing.imported-weyl-scalars-petrov-type-schwarzschild}.

A blow-up seen in one gauge might still be an artefact of that gauge.

This step closes the loophole with the curvature scalar \[ K = R\_{abcd}R\^{}{abcd} = \frac{48\,M\^{}2}{r\^{}6} \] which every gauge transformation leaves unchanged and which diverges as $r \to 0$. The division-by-zero locus is therefore a property of the solution, not of its presentation.

\begin{proposition}[Non-removability: the invariant argument]\label{c:singularity.non-removability-invariant-argument}
At the centre of a non-rotating black hole, a measure of curvature that does not depend on how the description is set up grows without bound. Because every change of description allowed in the theory leaves this quantity unchanged, no choice of description can make the blow-up disappear. The breakdown at the centre is therefore a feature of the solution itself rather than an artefact of how it is written down.
\emph{First proved in} T5, Prop 2.7 (\href{https://doi.org/10.5281/zenodo.22721790}{doi:10.5281/zenodo.22721790}).
\end{proposition}
\noindent\emph{Rests on:} \ref{c:singularity.solution-where-divergence-lies}.

Knowing that something genuinely fails, the next question is what shape the failure takes.

Divided by its growth, the field tends to a rank-one limit, $\beta^{-1}\hb \to -P$, and the metric to \[ \hat g(a,b) = -(a\cdot\et)(b\cdot\et) \] a form of signature $(0,1,3)$ that sees only the time direction. The radial direction $\er$ joins its radical only in the limit, never at any finite radius.

\begin{proposition}[The degeneration, exactly]\label{c:singularity.degeneration-exactly}
Approaching the centre of a black hole the gravitational field grows without bound, but once that overall growth is divided out the field settles to a definite limit. The limit is degenerate: it registers only the time direction and is blind to the other three. The radial direction, along which the light cones fold, becomes one of these blind directions only in the limit, not at any finite distance from the centre.
\emph{First proved in} T5, Prop 2.6 (\href{https://doi.org/10.5281/zenodo.22721790}{doi:10.5281/zenodo.22721790}).
\end{proposition}

The same collapse shows up in the light cones.

The two radial null directions are $\et + (-\beta \pm 1)\er$, and as $r \to 0$ both tip inward and merge onto the single direction $\er$, while $g(\et,\et)$ diverges. No gauge choice removes this degeneration through a null kernel direction.

\begin{proposition}[The cones fold]\label{c:singularity.cones-fold}
Light cones record the directions in which light can travel at each point. Following the two radial light directions inward towards the centre of a black hole, both swing round to point inward and, in the limit at the centre, merge into a single direction. At the same time the time part of the metric grows without bound, so the classical field cannot be rescaled to stay finite and well behaved there, and the failure cannot be removed by a change of gauge.
\emph{First proved in} T5, Prop 2.5 (\href{https://doi.org/10.5281/zenodo.22721790}{doi:10.5281/zenodo.22721790}).
\end{proposition}
\noindent\emph{Rests on:} \ref{c:singularity.degeneration-exactly}, \ref{c:singularity.non-removability-invariant-argument}, \ref{c:std-algebra.orthogonal-anticommute}.

The series classifies the ways in which division can fail, and this step places the black-hole centre in that classification.

The obstruction is of type N, a null direction with no metric partner. Type C cannot occur, because $\Cl(1,3) \cong \Mat_2(\HH)$ is central simple, and type P is present but imposes no transport.

That diagnosis decides the repair: supply the missing partner by the Witt completion.

\begin{proposition}[Type N, and only type N]\label{c:singularity.type-n-only-type-n}
At the centre of a non-rotating black hole one direction of the geometry collapses, and this result identifies which kind of algebraic obstruction that is in the series' classification. It is the kind caused by a direction with no partner to pair with. The other substantive kind, which would need the algebra to split into independent pieces, cannot occur for the algebra of spacetime, and a third kind is present but harmless. So the repair that applies is the one that supplies the missing partner by enlarging the algebra.
\emph{First proved in} T5, Prop 3.1 (\href{https://doi.org/10.5281/zenodo.22721790}{doi:10.5281/zenodo.22721790}).
\end{proposition}
\noindent\emph{Rests on:} \ref{c:sheaf.localised-transport}, \ref{c:sheaf.one-sided-completion}, \ref{c:sheaf.one-sided-completion-singularity}, \ref{c:sheaf.schwarzschild-degeneration-boundary-tame}, \ref{c:singularity.degeneration-exactly}, \ref{c:tower.rigidity-clock-frozen}, \ref{c:transport.existence-classification-type-c}, \ref{c:transport.existence-locus-type-n}, \ref{c:transport.type-p-self-repairing}, \ref{c:primer.classification-real-clifford-algebras}, \ref{c:std-algebra.sta-matrix}.

The repair is carried over from the series' general treatment of degenerating geometries.

A hyperbolic plane $(\nO, \nInf)$ is adjoined and coupled to the dying direction, with $\mu\nu < 0$ and a bump inside the horizon, and the ambient form then has constant signature $(2,4)$ everywhere, the singularity included. The ladder pair that partners the lost direction exists exactly over the singularity, and no shadow wall forms anywhere.

\begin{assembled}[Imported: one-sided completion at the singularity]\label{c:singularity.imported-one-sided-completion-singularity}
The centre of a non-rotating black hole is where the geometry loses a direction. This result, carried over from the sheaf paper, adds two extra directions and couples them to the lost one, with couplings of opposite sign switched on only inside the horizon, so that the enlarged geometry is non-degenerate everywhere, the centre included. The new partner for the lost direction exists exactly at the centre, and the repair creates no new defect anywhere else.
\emph{First proved in} T5, Fact B.2 (\href{https://doi.org/10.5281/zenodo.22721790}{doi:10.5281/zenodo.22721790}).
\end{assembled}
\noindent\emph{Rests on:} \ref{c:sheaf.one-sided-completion-singularity}, \ref{c:sheaf.schwarzschild-degeneration-boundary-tame}.

\section{Results used from elsewhere}
The proofs above use the following results of the series, stated briefly; each is proved in the paper named.

\begin{fact}[Imported: the Weyl scalars, and the Petrov type of Schwarzschild]\label{c:pairing.imported-weyl-scalars-petrov-type-schwarzschild}
The Weyl map of the field strength has five complex null-tetrad components $\Psi_0, \dots, \Psi_4$, and the Schwarzschild field is of Petrov type D with only $\Psi_2 = -M/r^3$ nonzero, its repeated principal directions radial.
\emph{Source:} R. Penrose, W. Rindler, Spinors and Space-Time, Vol. 2: Spinor and Twistor Methods in Space-Time Geometry, Cambridge University Press (1986)., Ch.~8; E. Newman, R. Penrose, An approach to gravitational radiation by a method of spin coefficients, J. Math. Phys. 3 (1962) 566--578..
\end{fact}

\begin{fact}[Classification of the real Clifford algebras]\label{c:primer.classification-real-clifford-algebras}
$\Cl(p,q)$ is a matrix algebra over $\R$, $\C$ or $\HH$, or a sum of two such, according to $p - q \bmod 8$, and it is simple except when $p - q \equiv 1 \pmod 4$.
\emph{Source:} T0, Thm 9.3 (\href{https://doi.org/10.5281/zenodo.22714501}{doi:10.5281/zenodo.22714501}).
\end{fact}

\begin{fact}[Localised transport]\label{c:sheaf.localised-transport}
The degenerate Clifford sheaf embeds canonically in the completed sheaf, and on the wall the radical covector acquires a partner with which it forms a ladder pair that exists exactly over the wall.
\emph{Source:} T4, Cor 4.5 (\href{https://doi.org/10.5281/zenodo.22721669}{doi:10.5281/zenodo.22721669}).
\end{fact}

\begin{fact}[One-sided completion]\label{c:sheaf.one-sided-completion}
Under tameness, with signature $(p,q)$, $p \geq 1$, on the physical side and a trivial radical line bundle, coupling an adjoined hyperbolic plane to the dying direction gives an ambient form of constant signature $(p{+}1, q{+}1)$ on the whole physical side, wall included.
\emph{Source:} T4, Thm 4.4 (\href{https://doi.org/10.5281/zenodo.22721669}{doi:10.5281/zenodo.22721669}).
\end{fact}

\begin{fact}[One-sided completion at the singularity]\label{c:sheaf.one-sided-completion-singularity}
At the blown-up Schwarzschild singularity, coupling an adjoined hyperbolic plane to the dying direction with $\mu\nu < 0$ and bump inside the horizon gives an ambient form of constant signature $(2,4)$ everywhere, boundary included, with no shadow wall and a ladder pair existing exactly over the boundary.
\emph{Source:} T4, Thm 8.12 (\href{https://doi.org/10.5281/zenodo.22721669}{doi:10.5281/zenodo.22721669}).
\end{fact}

\begin{fact}[The Schwarzschild degeneration is boundary-tame in $\shW$]\label{c:sheaf.schwarzschild-degeneration-boundary-tame}
In the weighted module at the blown-up Schwarzschild singularity, the derived cometric of the Painlev\'e-Gullstrand gauge field extends smoothly to the boundary and degenerates there with corank exactly one and to exactly first order, with radical $k = \rho + m\,\dtPG$ annihilating the infalling congruence.
\emph{Source:} T4, Thm 8.7 (\href{https://doi.org/10.5281/zenodo.22721669}{doi:10.5281/zenodo.22721669}).
\end{fact}

\begin{fact}[Imported: equivalence with the Einstein equations]\label{c:singularity.imported-equivalence-einstein-equations}
Where $\hb$ is invertible, gauge-theory gravity is equivalent to the Einstein equations and its curvature bivector reproduces the Riemann tensor.
\emph{Source:} A. Lasenby, C. Doran and S. Gull, Gravity, gauge theories and geometric algebra, Phil. Trans. R. Soc. Lond. A 356 (1998), 487--582., \S 4 (\S 4.3, covariant forms of the field equations), App.~C; C. Doran and A. Lasenby, Geometric Algebra for Physicists, Cambridge University Press (2003)., Ch.~13-14.
\end{fact}

\begin{fact}[Orthogonal vectors anticommute]\label{c:std-algebra.orthogonal-anticommute}
$uv = -vu$ iff $B(u,v) = 0$; in an orthogonal basis $e_I e_J = \pm \prod_{I \cap J} Q(e_i)\; e_{I \triangle J}$ with the sorting sign.
\emph{Source:} the algebra core of the Transport Lab, proved there with an exact check; standard, see C. Doran and A. Lasenby, Geometric Algebra for Physicists, Cambridge University Press (2003), Ch. 4 (foundations of geometric algebra).
\end{fact}

\begin{fact}[The spacetime algebra as a matrix algebra, and the spacetime split]\label{c:std-algebra.sta-matrix}
$\Cl(1,3) \cong \Mat_2(\HH)$, the Dirac matrices represent $\C \otimes \Cl(1,3) \cong \Mat_4(\C)$, and $\sigma_k = \gamma_k\gamma_0$ give $\Cl^{+}(1,3) \cong \Cl(3,0)$ with $\sigma_1\sigma_2\sigma_3 = I$.
\emph{Source:} the algebra core of the Transport Lab, proved there with an exact check; standard, see C. Doran and A. Lasenby, Geometric Algebra for Physicists, Cambridge University Press (2003), Ch. 5 (relativity and spacetime).
\end{fact}

\begin{fact}[Rigidity; the clock is frozen]\label{c:tower.rigidity-clock-frozen}
Every floor of the tower shares the ground floor's charge, division-ring type, centre, simplicity class, Morita class and obstruction taxonomy, so the physical tower $\Cl(1{+}k, 3{+}k) \cong \Mat_{2^{k+1}}(\HH)$ is frozen at clock position $6$.
\emph{Source:} T3, Thm 3.6 (\href{https://doi.org/10.5281/zenodo.22721556}{doi:10.5281/zenodo.22721556}).
\end{fact}

\begin{fact}[Existence and classification of type C]\label{c:transport.existence-classification-type-c}
$\Cl(p,q)$, $n \geq 1$, has a nontrivial central idempotent exactly when $p - q \equiv 1 \pmod 4$, and then the only ones are $\varepsilon_\pm = \tfrac12(1 \pm \omega)$, which split the algebra into two simple ideals and are exchanged by the grade involution but by no inner automorphism.
\emph{Source:} T1, Thm 3.6 (\href{https://doi.org/10.5281/zenodo.21743863}{doi:10.5281/zenodo.21743863}).
\end{fact}

\begin{fact}[Existence locus of type N]\label{c:transport.existence-locus-type-n}
$\Cl(p,q)$ contains null vectors exactly when $pq \neq 0$; they then form a cone of dimension $n - 1$, and each is a two-sided zero divisor with no left or right inverse in any unital algebra containing $\Cl(p,q)$.
\emph{Source:} T1, Prop 3.3 (\href{https://doi.org/10.5281/zenodo.21743863}{doi:10.5281/zenodo.21743863}).
\end{fact}

\begin{fact}[Type P is self-repairing]\label{c:transport.type-p-self-repairing}
In any unital ring the identity assignment $\sigma(e) = e$ on idempotents satisfies both von Neumann identities and is equivariant under every automorphism of every extension preserving the ring, so type P obstructions are repaired at signature increment $(0,0)$.
\emph{Source:} T1, Prop 6.1 (\href{https://doi.org/10.5281/zenodo.21743863}{doi:10.5281/zenodo.21743863}).
\end{fact}

\section*{Conclusion}
The thread runs from a field that inverts everywhere away from the centre, through a divergence no gauge choice removes, to a degenerate limit with one partnerless null direction and the completion that supplies its partner.

This walk adds no postulate of its own; it rests on imported results, above all the equivalence of Gauge Theory Gravity with Einstein's equations wherever $\hb$ is invertible. The repair uses the general completion only after the black-hole degeneration has been shown to be of the tame kind the completion requires.

\end{document}